% pandoc-defs.tex -- pandoc-friendly redefinitions of Trench's custom macros for
% the HTML/EPUB web edition (impo trench family). Prepended to the (\over-converted)
% source so pandoc EXPANDS these to constructs it understands; the PDF's real shim
% (trench-osbook.sty) is untouched. Spike version -- grows as chapters surface more.

% --- sectioning ---
% NOTE: sectioning is rewritten to real \chapter/\section as TEXT by the web
% preprocessor, NOT via macros here -- pandoc empties the whole document if two
% \newcommands both expand to sectioning commands (\chapter and \section). So the
% preprocessor does:  \chaptertitle{X}->\chapter{X}, \newsection{a}{b}{c}->\section{c},
% \sectiontitle{X}->(removed). Only these harmless leftovers are defined here:
% (\thissection is \newcommand'd by the body itself; do not predefine it here.)
% PDF-bookmark commands (hyperref/bookmark) -> no-ops for the web edition.
\providecommand{\currentpdfbookmark}[2]{}
\providecommand{\subpdfbookmark}[2]{}
% counters the body references (exercise/place; list counters lcal/rm)
\newcounter{exercise}\newcounter{place}\newcounter{lcal}\newcounter{rm}
% the answer key does \renewcommand{\newsection}[1]{...}; give it a dummy to renew
\providecommand{\newsection}{}
\newcommand{\markonleft}[1]{}
\newcommand{\markonright}[1]{}

% --- theorem-likes -> simple labelled blocks ---
\newenvironment{theorem}{\par\medskip\noindent\textbf{Theorem.}\itshape\ }{\par\medskip}
\newenvironment{example}{\par\medskip\noindent\textbf{Example.}\ }{\par\medskip}
% NOTE: no `\ ` after #1 -- the source writes `\begin{definition}\color{blue}...`
% (unbraced), so #1 grabs \color and a trailing `\ ` would become \color's arg.
\newenvironment{definition}[1]{\par\medskip\noindent\textbf{Definition.}\ #1}{\par\medskip}
\newcommand{\proof}{\par\noindent\textit{Proof.}\ }
\def\endproof{}%  \newcommand refuses \end... names
\newcommand{\bbox}{}
\newcommand{\solution}{\par\noindent\textit{Solution.}\ }
\newcommand{\solutionpart}[1]{\par\noindent\textit{Solution (#1).}\ }
\newcommand{\remark}[1]{\par\noindent\textbf{Remark.}\ #1}

% --- exercises: page-refs become HYPERLINKS to the target (maintainer-approved) ---
\newcommand{\exercises}{\par\medskip\noindent\textbf{Exercises}\par}
\newenvironment{exerciselist}{\begin{enumerate}}{\end{enumerate}}
\newcommand{\exercisetext}[1]{\end{enumerate}#1\begin{enumerate}}
\newcommand{\exer}[1]{\textbf{#1.}\ }
% \answer{N} used to print "N (p. <page>)"; now: link to the exercise, no page number.
\newcommand{\answer}[1]{\ (\ref{exer:#1})}
\renewcommand{\pageref}[1]{\ref{#1}}   % any surviving \pageref -> a link, not a page
\newcommand{\Cex}{\textbf{[C]}\ }
\newcommand{\CGex}{\textbf{[C/G]}\ }
\newcommand{\Lex}{\textbf{[L]}\ }
\newcommand{\technology}{\par\noindent\textbf{USING TECHNOLOGY}\par}
\newenvironment{alist}{\begin{enumerate}}{\end{enumerate}}
\newenvironment{Alist}{\begin{enumerate}}{\end{enumerate}}
\newenvironment{rmlist}{\begin{enumerate}}{\end{enumerate}}
\newenvironment{numberlist}{\begin{enumerate}}{\end{enumerate}}
\newenvironment{newtable}[1]{\begin{quote}\textbf{Table.}\ #1}{\end{quote}}

% --- boxes / spacing helpers -> passthrough ---
\newcommand{\place}{\par}
\newcommand{\boxit}[1]{#1}
\newcommand{\helpbox}{}
\newcommand{\bigbox}[1]{#1}
\newcommand{\blankbox}[1]{#1}
\newcommand{\topfig}[1]{}
\newcommand{\hint}[1]{\textsc{Hint}: \textit{#1}}
\newcommand{\note}[1]{(\textsc{Note}: #1)}
\renewcommand{\part}[1]{(\textbf{#1})}

% --- old LaTeX 2.09 font declarations (Trench uses \bf \it \rm \sc in text AND math) ---
\DeclareOldFontCommand{\bf}{\normalfont\bfseries}{\mathbf}
\DeclareOldFontCommand{\it}{\normalfont\itshape}{\mathit}
\DeclareOldFontCommand{\rm}{\normalfont\rmfamily}{\mathrm}
\DeclareOldFontCommand{\sf}{\normalfont\sffamily}{\mathsf}
\DeclareOldFontCommand{\tt}{\normalfont\ttfamily}{\mathtt}
\DeclareOldFontCommand{\sc}{\normalfont\scshape}{\mathrm}
\DeclareOldFontCommand{\sl}{\normalfont\itshape}{\mathit}
\DeclareOldFontCommand{\cal}{\normalfont}{\mathcal}

% --- math shorthands ---
\newcommand{\dst}{\displaystyle}
\newcommand{\R}{\mathbb{R}}
\newcommand{\E}{\mathbb{E}}
\newcommand{\set}[2]{\left\{#1 \,\middle|\, #2\right\}}
\newcommand{\dist}{\operatorname{dist}}
\newcommand{\adj}{\operatorname{adj}}
\newcommand{\df}{=_{\mathrm{df}}}
\newcommand{\lims}[2]{\Big|_{#1}^{#2}}
\newcommand{\eqline}[1]{#1}
\newcommand{\twocol}[2]{\left[\begin{array}{l}#1\\#2\end{array}\right]}
\newcommand{\ctwocol}[2]{\left[\begin{array}{c}#1\\#2\end{array}\right]}
\newcommand{\threecol}[3]{\left[\begin{array}{r}#1\\#2\\#3\end{array}\right]}
\newcommand{\cthreecol}[3]{\left[\begin{array}{c}#1\\#2\\#3\end{array}\right]}
\newcommand{\twobytwo}[4]{\left[\begin{array}{rr}#1&#2\\#3&#4\end{array}\right]}
\newcommand{\threebythree}[9]{\left[\begin{array}{rrr}#1&#2&#3\\#4&#5&#6\\#7&#8&#9\end{array}\right]}
% more matrix/array builders used in the later chapters (from the shim)
\newcommand{\col}[2]{\left[\begin{array}{c}#1_1\\#1_2\\\vdots\\#1_#2\end{array}\right]}
\newcommand{\colfunc}[2]{\left[\begin{array}{c}#1_1(t)\\#1_2(t)\\\vdots\\#1_#2(t)\end{array}\right]}
\newcommand{\cthreebythree}[9]{\left[\begin{array}{ccc}#1&#2&#3\\#4&#5&#6\\#7&#8&#9\end{array}\right]}
\newcommand{\twochar}[4]{\left|\begin{array}{cc}#1-\lambda&#2\\#3&#4-\lambda\end{array}\right|}
\newcommand{\threechar}[9]{\left[\begin{array}{ccc}#1-\lambda&#2&#3\\#4&#5-\lambda&#6\\#7&#8&#9-\lambda\end{array}\right]}
\newcommand{\matfunc}[3]{\left[\begin{array}{cccc}#1_{11}(t)&#1_{12}(t)&\cdots&#1_{1#3}(t)\\%
  #1_{21}(t)&#1_{22}(t)&\cdots&#1_{2#3}(t)\\\vdots&\vdots&\ddots&\vdots\\%
  #1_{#21}(t)&#1_{#22}(t)&\cdots&#1_{#2#3}(t)\end{array}\right]}
\newcommand{\arraytext}[1]{\noalign{\hspace*{-1.75em}\mbox{#1}\vspace*{-9\jot}}\nonumber\\}
